\section{总结}
本文围绕“量子信息如何合理进入经典同化器”这一核心问题，给出了一组基于 Lorenz96 的量子弱融合受控实验。实验表明，量子结构本身并非无效，但若直接主导观测权重或分析更新，容易破坏经典同化的稳定性；相反，当其被解释为局地弱几何修正信号时，可以在保持 LETKF 稳定骨架的同时带来温和但可量化的收益。

因此，本文给出的主要结论不是“量子模块替代了经典同化器”，而是“在当前受控设定下，量子模块在局地结构建模和弱融合层面显示出一定现实价值”。从最终验证结果看，最优量子弱融合方案在 test\_1 上优于 fixed 基线，同时避免了 corr 方案的明显退化，说明“稳定骨架 + 局地弱修正”值得作为后续量子增强数据同化研究的一条候选方向。需要强调的是，本文的证据规模目前仍局限于单一 Lorenz96 设定、有限参数网格和单测试集比较，后续工作将进一步围绕更完整的验证划分、更丰富的中间消融、更稳定的量子表示层以及 gated/adaptive 变体的统一比较展开。
